\documentclass[10pt,a4paper]{amsart}
\usepackage[margin=19mm]{geometry}
\usepackage{amsmath,amssymb,mathtools}
\usepackage[scr=rsfs,cal=euler]{mathalfa}
\usepackage{xfrac}
\usepackage[unicode,hidelinks,bookmarks=false]{hyperref}
\newcommand{\ie}{i.e.\ }
\newcommand{\cf}{cf.\ }
\newcommand{\resp}{respectively\ }
\newcommand{\Subsection}{\subsection}
\providecommand{\nobreakdash}{\nobreak\hbox{-}}
\setlength{\emergencystretch}{2em}
\multlinegap0pt
\usepackage{bm}
\usepackage{bbm}
\usepackage{dsfont}
\usepackage{xspace}
\usepackage{cancel}
\usepackage{mathtools}
\usepackage{amsthm}
\makeatletter
\@ifundefined{theorem}{\newtheorem{theorem}{Theorem}}{}
\@ifundefined{definition}{\newtheorem{definition}[theorem]{Definition}}{}
\@ifundefined{prop}{\newtheorem{prop}[theorem]{Proposition}}{}
\makeatother
\DeclareMathOperator{\Mod}{Mod}
\DeclareMathOperator{\PGL}{PGL}
\DeclareMathOperator{\Sym}{Sym}
\newcommand{\C}{\mathcal{C}}
\newcommand{\M}{\mathcal{M}}
\newcommand{\QQ}{\mathbb{Q}}
\newcommand{\ZZ}{\mathbb{Z}}
\newcommand{\mc}{\mathcal}
\newcommand{\tW}{\widetilde{W}}
\title[Stable cohomology of universal character varieties]{Stable cohomology of universal character varieties}
\author{Ishan Banerjee, Faye Jackson, Anne Larsen, Sam Payne,, Xiyan Zhong}
\date{Source: arXiv:2512.13879v3 (CC BY 4.0)}
\begin{document}
\maketitle

\noindent\emph{Source and licence.} Stable cohomology of universal character varieties. Version arXiv:2512.13879v3 (CC BY 4.0), distributed under the Creative Commons Attribution 4.0 International licence, as recorded in the arXiv licence metadata for this exact version and shown on the abstract page https://arxiv.org/abs/2512.13879v3. Full rights notice: licenses/arxiv-2512.13879-CC-BY-4.0.txt.\par\smallskip

\noindent\textit{Editorial scope.} Counted unit: the Prop whose source label is \texttt{prop: monodromy action} in the article (source0.tex lines 911--916, source label \texttt{prop: monodromy action}), delivered together with the article's own complete original proof of it (source0.tex lines 918--946, one \texttt{proof} environment of 29 lines). No mathematics is written, repaired or re-derived: statement and proof are the article's own text line for line. Required context retained uncounted: defn: relative total PGL (lines 398--403); def: alpha, beta, psi (lines 871--873); prop: free generation (lines 877--879); prop: identifying monodromy (lines 887--889). Every definition, display and label of those lines is retained unchanged. Omitted from this package and not redistributed: 1--397, 404--870, 874--876, 880--886, 890--910, 917--917, 947--1792. Nothing inside a retained excerpt is omitted or altered: every retained body is delivered exactly as the article's lines are written, so no omission receipt of any kind accompanies this package. Citations: the retained excerpts carry no citation command at all, so no bibliography entry of the article is transcribed and no reference is redistributed. Reference adaptation: none was needed and none was made; every reference inside the delivered unit resolves inside it, so no unresolved reference or citation is produced. Printed slips kept literal: no printed word, symbol, hypothesis or number of the counted statement or of its proof was altered. Layout and class: the article's own class is \texttt{amsart} with packages including amsmath, amssymb, amsthm, mathrsfs, mathtools, setspace, hyperref, cleveref, xcolor, comment, tikz, tikz-cd, xy; the wrapper is the lane's pinned \texttt{amsart} editorial wrapper at 10pt on A4, which restates the theorem-like environments the retained text needs and adds the article's own macro definitions that the retained text needs (\texttt{\textbackslash C}, \texttt{\textbackslash M}, \texttt{\textbackslash Mod}, \texttt{\textbackslash PGL}, \texttt{\textbackslash QQ}, \texttt{\textbackslash Sym}, \texttt{\textbackslash ZZ}, \texttt{\textbackslash mc}, \texttt{\textbackslash tW}), with extra wrapper packages (bm, bbm, dsfont, xspace, cancel, mathtools). The delivered numbering is the unit's own. Assets: no retained span contains a figure environment, a graphics-inclusion command, a PDF-page inclusion, a table or a TikZ picture; no image file is copied into this package, the package declares no asset dependency and no image file is redistributed with it. Licence: version arXiv:2512.13879v3 of this article is distributed under the Creative Commons Attribution 4.0 International licence, as recorded in the arXiv licence metadata for this exact version and shown on the abstract page https://arxiv.org/abs/2512.13879v3. A full rights notice is delivered at licenses/arxiv-2512.13879-CC-BY-4.0.txt. Characters: every retained excerpt is pure ASCII, so the delivered engine, XeLaTeX with its default Latin Modern text font, reports no missing character.
\begin{definition}\label{defn: relative total PGL}
	For $d$ coprime to $n$, the \emph{universal degree $d$ $\PGL$-character variety} is 
	\begin{equation*}
		\M_{B}^{\PGL}(\C_g, n, d) \coloneqq \mc T_g \times_{\Mod_g} \M_{B}^{\PGL}(S_g,n,d).
	\end{equation*}
\end{definition}

\begin{definition}\label{def: alpha, beta, psi}
    For $i \in \{0,1,2\}$ and $j \in \ZZ_{\ge 2}$ we define the K\"unneth component $$V_{i,j} := \{p_{2*}(c_j(\mc P) \cup p_1^*(x)): x \in H^i(S_g;\QQ)\} \subset H^{i+2j-2}(\M_B^{\PGL}(S_g,n,d); \QQ).$$
\end{definition}

\begin{prop}\label{prop: free generation} %
    The subspaces $V_{i,j}$ for $0 \le i \le 2$ and $2 \le j \le n$ generate the ring $H^*(\M^{\PGL}_B(S_g,n,d); \QQ)$ with no relations in degrees less than or equal to $R_{g,n}$. Moreover, $\dim V_{i,j} = \dim H^i(S_g; \QQ)$.
\end{prop}

\begin{prop} \label{prop: identifying monodromy}
The Chern classes $c_i(\mc P) \in H^{2i}(S_g \times \M_B^{\PGL}(S_g,n,d); \QQ)$ are fixed by the mapping class group action.
\end{prop}

\begin{prop} \label{prop: monodromy action}
Let $\pi \colon \M_B^{\PGL}(\C_g,n,d) \to \M_g$ be the projection. In degrees less than or equal to $R_{g,n}$, there is an isomorphism of graded local systems on $\M_g$
\[
R^\bullet \pi_* \QQ \;\cong\; \textstyle{\bigwedge^\bullet}(V_g \otimes \tW_n) \otimes \Sym^\bullet(V_\circ \otimes \tW_n).
\]
\end{prop}

\begin{proof}
Since $\M_B^{\PGL}(\C_g,n,d)\to\M_g$ is the bundle associated to the $\Mod_g$-action on the fiber $\M_B^{\PGL}(S_g,n,d)$ (Definition~\ref{defn: relative total PGL}), the local system $R^q\pi_*\QQ$ is the one associated to the graded $\Mod_g$-module $H^q(\M_B^{\PGL}(S_g,n,d);\QQ)$. It therefore suffices to identify this module in degrees at most $R_{g,n}$.
For $0\le i\le 2$ and $2\le j\le n$, Definition~\ref{def: alpha, beta, psi} gives a surjection
\[
H^i(S_g;\QQ)\longrightarrow V_{i,j},
\qquad
x\longmapsto p_{2*}(c_j(\mc P)\cup p_1^*x).
\]
It is $\Mod_g$-equivariant by Proposition~\ref{prop: identifying monodromy}, and is an isomorphism by Proposition~\ref{prop: free generation}. Hence, as graded $\Mod_g$-modules,
\[
V_{0,j}\cong \QQ[2-2j],\qquad
V_{1,j}\cong V_g[2-2j],\qquad
V_{2,j}\cong \QQ[-2j].
\]
Equivalently, as graded $\Mod_g$-modules, we have
\[
\bigoplus_{j=2}^n V_{1,j}\cong V_g\otimes \tW_n,
\qquad
\bigoplus_{j=2}^n(V_{0,j}\oplus V_{2,j})\cong V_\circ\otimes \tW_n.
\]
The first summand consists of odd-degree generators, and the second of even-degree generators. Proposition~\ref{prop: free generation} therefore gives
\[
H^*(\M_B^{\PGL}(S_g,n,d);\QQ)\cong
\textstyle{\bigwedge^\bullet}(V_g\otimes \tW_n)
\otimes
\Sym^\bullet(V_\circ\otimes \tW_n)
\]
in degrees less than or equal to $R_{g,n}$, as required. 
\end{proof}

\end{document}
